\documentclass[11pt]{article}
\usepackage{fontspec}
\setmainfont{Times New Roman}
\setsansfont{Arial}
\setmonofont{Consolas}
\usepackage{amsmath,amssymb}
\usepackage{geometry}
\usepackage{microtype}
\geometry{a4paper,margin=3cm}
\setlength{\parindent}{1.25em}
\setlength{\parskip}{0pt}
\emergencystretch=4em
\allowdisplaybreaks
\let\A\relax
\newcommand{\A}{\mathfrak A}
\newcommand{\R}{\mathfrak R}

\begin{document}

\begin{center}
\textbf{\S\ 20. \quad Формы 7. поряда (систем $s_{III}$).}
\end{center}

Свијанье $s$ с $f\cdot \nu$, $a_{\nu}$, $a_{\nu}^2$.

\noindent\textit{Незводиме формы:}
\[
\begin{gathered}
s_{\nu}(asu),\quad a_{\nu}(asu),\quad s_{\nu}(asu)^2,\quad a_{\nu}(asu)^2,\quad
s_{\nu}(asu)^3,\quad a_{\nu}(asu)^3,\\
a_{\nu} s_{\nu},\qquad a_{\nu} s_{\nu}(asu),\qquad a_{\nu}^2(asu),\qquad a_{\nu}^2(asu)^2.
\end{gathered}
\]

\noindent\textit{Редукције:}

Само-разумльиве редукције, кторе наставајут из прјамого свијаньа с редуцентами $s_{\nu}^2$ и $s_\eta(Hsu)$, не будут посебно поменьене; такоже не разклад трансвекциј с функционалными детерминантами.

Форма $a_{\nu}s_{\nu}(asu)^2$ и $a_{\nu}^2s_{\nu}(asu)^2$: гл. формулу (19.), \S\ 12;
\[
a_{\nu}s_{\nu}(asu)^3=a_{\nu}(a\widehat{\nu}_1\nu_2u)^3(\nu_1\nu_2\nu)=0.
\]
Ряд форм $a_{\nu}^2s_{\nu}$: редуцент $s_\vartheta^2(\theta su)$ доставаје се повратом $a_{\nu}^2$ к нижшим формам, т. ј. двојној редукције изразов
\[
s_\vartheta^2(a\theta s)(a\theta u),\qquad
s_\vartheta^2(a\theta s)(a\theta u)^2
\]
по следујучем наставку:
\[
\begin{aligned}
s_\vartheta^2(a\theta s)(a\theta u)
&=[(a\theta u)^2]s^3+\frac15[a_\vartheta(a\theta u)^2]s^2+\frac1{60}[a_\vartheta^2(a\theta u)^2]s\\
&=\frac12a_{\nu}^2s_{\nu}-\frac23a_{\nu}s_{\nu}^2,\\
s_\vartheta^2(a\theta s)(a\theta u)^2
&=\frac45[a_\vartheta(a\theta u)^2]_{\widehat{us}}s^2+\frac{23}{180}[a_\vartheta^2(a\theta u)^2]_{\widehat{us}}s\\
&=-\frac12a_{\nu}^2s_{\nu}(asu).
\end{aligned}
\]
При том вельа $a_{\nu}^2b_{\nu}(abu)=0$.

Наставајут формулы:
\[
\tag{25}
\begin{array}{rcl@{\qquad}rcl}
\underline{a_{\nu}^2s_{\nu}}&=&\dfrac43i\cdot a_{\nu}^2b_{\nu},
& \underline{a_{\nu}s_{\nu}(asu)^2}&=&\dfrac13i\cdot\theta_{\nu}^2-\dfrac1{18}iH,\\[0.7em]
\underline{a_{\nu}s_{\nu}(asu)^3}&=&0,
& \underline{a_{\nu}^2s_{\nu}(asu)}&=&0,\\[0.7em]
\underline{a_{\nu}^2s_{\nu}(asu)^2}&=&0.
\end{array}
\]

\noindent\textit{Последки:}

1) $a_{\nu}s_{\nu}(asu)^2$, $a_{\nu}^2s_{\nu}$ сут редуценты.

2) Из того доставајут се вредности форм $(\Delta su)^3$, $(\Delta su)^4$, признаных редуцибилными в \S\ 19; подобно за соответны ряд форм $(agu)^2$, кторы при свијаньи с $\nu$ чрез редуцент $g_{\nu}$ даваје двојну редукцију.
\[
\tag{26}
\begin{aligned}
\mathrm{a)}\quad (\Delta su)^3&=-\frac65a_{\nu}^2(asu)+3a_{\nu}s_{\nu}(asu)+2i\theta_{\nu}(\vartheta\nu x),\\
\mathrm{b)}\quad (\Delta su)^4&=-\frac25a_{\nu}^2(asu)^2+\frac43i\theta_{\nu}^2+\frac19iH.
\end{aligned}
\]

Из формул (24.) и (26.), мод $(i,J)$ (ср. стр. 71):
\[
\tag{27}
\begin{aligned}
\mathrm{a)}\quad (agu)^2&=\frac23(\Delta su)^2-\frac43a_{\nu}s_{\nu}+\frac35s\cdot a_{\nu}^2,\\
\mathrm{b)}\quad (agu)^3&=2a_{\nu}s_{\nu}(asu)-a_{\nu}^2(asu)^2,\\
\mathrm{c)}\quad (agu)^4&=-a_{\nu}^2(asu)^2.
\end{aligned}
\]

\end{document}

